% Persona primaria: Preventor de campo
\begingroup
\fontsize{7.5}{9.7}\selectfont
\ifdefined\pcolw\else\newlength{\pcolw}\fi
\ifdefined\pfootw\else\newlength{\pfootw}\fi
\definecolor{PInk}{RGB}{12,74,140}
\definecolor{PTeal}{RGB}{13,148,136}
\definecolor{PBlue}{RGB}{37,99,235}
\definecolor{PSlate}{RGB}{71,85,105}
\definecolor{PEmr}{RGB}{5,150,105}
\definecolor{PAmb}{RGB}{217,119,6}
\definecolor{PRed}{RGB}{220,38,38}
\definecolor{PVio}{RGB}{124,58,237}
\definecolor{PInd}{RGB}{67,56,202}
\definecolor{PSky}{RGB}{2,132,199}
\definecolor{PPlum}{RGB}{190,24,93}
\newcommand{\PDot}[1]{%
  \tikz[baseline=-0.42ex]\fill[#1] (0,0) circle (1.85pt);%
  \hspace{0.48em}%
}
\newcommand{\PTit}[1]{\makebox[\linewidth][c]{\sffamily\bfseries #1}}
\tikzset{
  pchip/.style={
    sharp corners,
    rounded corners=7.7pt,
    inner xsep=6.8pt,
    inner ysep=0pt,
    minimum height=15.4pt,
    text height=0.68em,
    text depth=0.22em,
    font=\sffamily\bfseries\scriptsize,
    draw=none,
    line width=0pt,
    outer sep=0pt,
    align=center,
    anchor=base,
  },
}
\newcommand{\PPills}[6]{%
  \begin{tikzpicture}[baseline=0.22em, rounded corners=0pt, sharp corners]
    \node[pchip, fill=#1, text=#2] (p1) {#3};
    \node[pchip, fill=#4, text=#5, base right=4.2pt of p1] (p2) {#6};
  \end{tikzpicture}%
}
\newcommand{\PCaja}[4][south]{%
  \noindent
  \begin{tikzpicture}[baseline=(current bounding box.#1)]
    \begin{scope}[shift={(0.4pt,-1.8pt)}]
      \node[
        fill=black,
        opacity=0.18,
        text opacity=0,
        rounded corners=13pt,
        inner xsep=8.4pt,
        inner ysep=5.8pt,
        outer sep=0pt,
        align=left,
        font=\fontsize{7.5}{9.7}\selectfont\sffamily,
        text width=\dimexpr#2-16.8pt\relax,
      ] {#4};
    \end{scope}
    \node[
      fill=#3!9,
      rounded corners=13pt,
      inner xsep=8.4pt,
      inner ysep=5.8pt,
      outer sep=0pt,
      text=black,
      align=left,
      font=\fontsize{7.5}{9.7}\selectfont\sffamily,
      text width=\dimexpr#2-16.8pt\relax,
    ] (pcard) {#4};
    \fill[#3, rounded corners=2.1pt]
      ([xshift=2.2pt,yshift=5.2pt]pcard.south west)
      rectangle
      ([xshift=5.7pt,yshift=-5.2pt]pcard.north west);
  \end{tikzpicture}%
}

\begin{tikzpicture}
  \node[
    rounded corners=16pt,
    inner sep=0pt,
    outer sep=0pt,
    fill=black!3.5,
  ] {%
    \begin{minipage}{0.988\linewidth}
      \begin{tikzpicture}
        \node[
          fill=PInk,
          rounded corners=12pt,
          inner sep=0pt,
          outer sep=0pt,
          text width=\linewidth-1.4pt,
        ] {%
          \begin{minipage}{\linewidth}
            \vspace{7pt}
            \hspace{9pt}%
            \begin{minipage}{0.60\linewidth}
              \begin{tikzpicture}[baseline=-0.55ex]
                \fill[white] (0,0) circle (0.34cm);
                \node[font=\sffamily\bfseries\scriptsize, color=PInk] at (0,0) {MR};
              \end{tikzpicture}%
              \hspace{6pt}%
              \raisebox{0.15ex}{\textcolor{white}{\sffamily\bfseries\Large Mart\'{\i}n R\'{\i}os}}\\[3pt]
              \hspace{1pt}\textcolor{white}{\sffamily\scriptsize Llegar a la visita con el historial, no con la lista de vencimientos}\\[2pt]
              \textcolor{white!88}{\sffamily\tiny Preventor de campo $\cdot$ ART $\cdot$ persona primaria}
            \end{minipage}%
            \hfill
            \begin{minipage}{0.34\linewidth}
              \raggedleft
              \PPills{white}{PInk}{Campo}{PSky}{white}{Escritorio y visita}\\[5pt]
              \textcolor{white!88}{\sffamily\tiny Uso diario, obligatorio}
            \end{minipage}%
            \hspace{9pt}%
            \vspace{7pt}
          \end{minipage}%
        };
      \end{tikzpicture}

      \vspace{8pt}
      \raggedright
      \begin{minipage}{\dimexpr\linewidth-22pt\relax}
        \setlength{\pcolw}{\dimexpr(\linewidth-16pt)/3\relax}
        \setlength{\pfootw}{\dimexpr(\linewidth-8pt)/2\relax}

        \PCaja{\linewidth}{PTeal}{%
          {\color{PTeal}\PTit{Qui\'en es y c\'omo trabaja}}\\[4pt]
          Recorre los establecimientos de su zona con una capacidad de visitas limitada.
          Arma la ronda con renovaci\'on de p\'oliza, frecuencia por actividad, RGRL/RAR y disparadores graves (ROAM).
          Coordina fecha con el empleador: si no hay agenda, la cobertura se atrasa.
          En el establecimiento emite informe para ART y empleador; el incumplimiento puede pegar en la al\'{\i}cuota y deja a la ART expuesta ante la SRT.
          Quiere llegar ya sabiendo qu\'e lesiones se reiteran. El score de Prevex queda interno: avisarlo al empleador puede \enquote{levantar la perdiz}.
        }

        \vspace{9pt}
        \noindent
        \begin{minipage}[t]{\pcolw}
          \PCaja[north]{\pcolw}{PBlue}{%
            {\color{PBlue}\PTit{Contexto de uso}}\\[4pt]
            \PDot{PBlue}D\'onde: PC de escritorio en la ART; la visita es en el establecimiento.\\[2.6pt]
            \PDot{PBlue}Cu\'ando: arma el trimestre y ajusta la agenda cada semana.\\[2.6pt]
            \PDot{PBlue}Por qu\'e: el trabajo lo obliga; no es un tablero opcional.\\[2.6pt]
            \PDot{PBlue}Competencia: experto en el circuito de visita, no en el modelo.
          }
        \end{minipage}%
        \hspace{8pt}%
        \begin{minipage}[t]{\pcolw}
          \PCaja[north]{\pcolw}{PSlate}{%
            {\color{PSlate}\PTit{Herramientas de hoy}}\\[4pt]
            \PDot{PSlate}Sistema legado: vencimientos y ROAM, sin alerta por correo.\\[2.6pt]
            \PDot{PSlate}Criterio propio sobre la lista.\\[2.6pt]
            \PDot{PSlate}Informe de visita y revisita.\\[2.6pt]
            \PDot{PSlate}Historial de accidentes del establecimiento, cuando lo consigue.
          }
        \end{minipage}%
        \hspace{8pt}%
        \begin{minipage}[t]{\pcolw}
          \PCaja[north]{\pcolw}{PEmr}{%
            {\color{PEmr}\PTit{Motivaciones}}\\[4pt]
            \PDot{PEmr}Cubrir la zona cr\'{\i}tica con el tiempo que tiene.\\[2.6pt]
            \PDot{PEmr}Salir de la visita con una recomendaci\'on fundada.\\[2.6pt]
            \PDot{PEmr}Dejar constancia del rol de la ART ante la SRT.\\[2.6pt]
            \PDot{PEmr}No improvisar in situ; el score no sale al empleador.
          }
        \end{minipage}

        \vspace{9pt}
        \noindent
        \begin{minipage}[t]{\pcolw}
          \PCaja[north]{\pcolw}{PAmb}{%
            {\color{PAmb}\PTit{Objetivos}}\\[4pt]
            \PDot{PAmb}Priorizar el trimestre, no la p\'oliza aislada.\\[2.6pt]
            \PDot{PAmb}Abrir el establecimiento y entender por qu\'e est\'a arriba.\\[2.6pt]
            \PDot{PAmb}Convertir el ranking en una visita agendada y cerrada.\\[2.6pt]
            \PDot{PAmb}Llevar historial de lesiones y patrones (p.\ ej.\ pinchazos en cosecha).
          }
        \end{minipage}%
        \hspace{8pt}%
        \begin{minipage}[t]{\pcolw}
          \PCaja[north]{\pcolw}{PRed}{%
            {\color{PRed}\PTit{Frustraciones}}\\[4pt]
            \PDot{PRed}No hay ranking predictivo: manda el calendario y el ROAM.\\[2.6pt]
            \PDot{PRed}Cuello de botella: su agenda + la del empleador.\\[2.6pt]
            \PDot{PRed}Listados sin notificaci\'on; todo depende de que \'el mire el sistema.\\[2.6pt]
            \PDot{PRed}Si el puntaje se filtra, alguien podr\'{\i}a aprovechar un punto d\'ebil.
          }
        \end{minipage}%
        \hspace{8pt}%
        \begin{minipage}[t]{\pcolw}
          \PCaja[north]{\pcolw}{PVio}{%
            {\color{PVio}\PTit{H\'abitos}}\\[4pt]
            \PDot{PVio}Combina el listado del sistema con lo que vio en visitas previas.\\[2.6pt]
            \PDot{PVio}Ciclo: visita $\to$ informe $\to$ recomendaciones $\to$ revisita.\\[2.6pt]
            \PDot{PVio}Documenta para la ART aun si el empleador no ejecuta.\\[2.6pt]
            \PDot{PVio}Desconf\'{\i}a de alertas abiertas hacia afuera del \'area.
          }
        \end{minipage}

        \vspace{10pt}
        \PCaja{\linewidth}{PInd}{%
          {\color{PInd}\PTit{Escenario}}\\[4pt]
          Arranca el trimestre, abre Prevex y ve la lista operativa de 1000 con bandas.
          Filtra su zona, entra a un establecimiento cr\'{\i}tico, lee SHAP y el resumen sin identificadores, y agenda la visita.
          En el predio ya sabe d\'onde insistir (protecci\'on visual, un puesto, un patr\'on de d\'{\i}as).
          Cierra el informe, marca la visita realizada y deja el score dentro del equipo.
        }

        \vspace{9pt}
        \noindent
        \begin{minipage}[t]{\pfootw}
          \PCaja[north]{\pfootw}{PSky}{%
            {\color{PSky}\PTit{Qu\'e necesita de Prevex}}\\[4pt]
            \PDot{PSky}Ranking trimestral top-1000 con bandas (HU-02, HU-03 / RF-02, RF-03, RF-12).\\[2.6pt]
            \PDot{PSky}Ficha con SHAP y resumen sin identificadores (HU-04 / RF-04, RNF-06, RNF-09).\\[2.6pt]
            \PDot{PSky}Agenda propia para programar y cerrar (HU-05, HU-06 / RF-05).\\[2.6pt]
            \PDot{PSky}El score no se exporta al empleador.
          }
        \end{minipage}%
        \hspace{8pt}%
        \begin{minipage}[t]{\pfootw}
          \PCaja[north]{\pfootw}{PPlum}{%
            {\color{PEmr}\PTit{\'Exito}}\\[3pt]
            La ronda cubre primero lo cr\'{\i}tico del trimestre y cada visita sale con una recomendaci\'on concreta.\\[5pt]
            {\color{PPlum}\PTit{Temor}}\\[3pt]
            No alcanzar la cobertura; que la SRT no encuentre registro; que alguien use el puntaje fuera de la ART.
          }
        \end{minipage}

        \vspace{4pt}
      \end{minipage}
    \end{minipage}%
  };
\end{tikzpicture}
\endgroup
